% answers_short.tex —— 2015 年·湖南卷·理科 · 简短答案速查
% 由 tools/make_short_answers.py 自 answers.tex 的 \gksolhead/\gkanswer 提取，
% 并与 2015/2015_数学真题_湖南理_答案要点.md 逐题比对后生成，勿手改（改源 MD 后重跑）。

\begin{gknotice}
\textbf{使用说明}：本节只给最终答案便于快速核对，完整推导见解析版。答案为编写者独立推导，\textbf{非官方标准答案}。
\end{gknotice}

\gkpart{选择题}
\begin{center}\begin{tabular}{ccccccccccc}
\toprule
题号 & 1 & 2 & 3 & 4 & 5 & 6 & 7 & 8 & 9 & 10 \\
\midrule
答案 & D & C & B & A & A & D & C & B & D & A \\
\bottomrule
\end{tabular}\end{center}

\gkpart{填空题}
\begin{itemize}[leftmargin=2.6em]
  \item[\textbf{11.}] \(0\)
  \item[\textbf{12.}] \(4\)
  \item[\textbf{13.}] \(\sqrt5\)
  \item[\textbf{14.}] 若按相邻项比值定义，则 \(a_n=3^{n-1}\)；若仅按递推式定义，则还包括 \(q=0\)、\((a_n)=(1,0,0,\ldots)\)。约定尚待来源核实，以上均为非官方条件结论。
  \item[\textbf{15.}] \((-\infty,0)\cup(1,+\infty)\)
\end{itemize}

\gkpart{解答题}
\begin{itemize}[leftmargin=2.6em]
  \item[\textbf{16.}] (1) A：\(\angle MEN+\angle NOM=180^\circ\)；B：\((x-1)^2+y^2=1\)；C：\(a+b\geqslant2\)；(2) A：\(FE\cdot FN=FM\cdot FO\)；B：\(|MA||MB|=18\)；C：两个不等式不能同时成立
  \item[\textbf{17.}] (1) \(B-A=\dfrac{\pi}{2}\)；(2) \((\sin A+\sin C)\in(\dfrac{\sqrt2}{2},\dfrac98]\)
  \item[\textbf{18.}] (1) \(\dfrac7{10}\)；(2) \(P(X=0)=\dfrac{64}{125},\ P(X=1)=\dfrac{48}{125},\ P(X=2)=\dfrac{12}{125},\ P(X=3)=\dfrac1{125};\ E(X)=\dfrac35\)
  \item[\textbf{19.}] (1) \(AB_1\perp PQ\)；(2) 四面体\(ADPQ\)的体积为\(24\)
  \item[\textbf{20.}] (1) \(\dfrac{y^2}{9}+\dfrac{x^2}{8}=1\)；(2) 直线斜率为\(\pm\dfrac{\sqrt6}{4}\)；题设同向条件下\(\angle MFD>90^\circ\)
  \item[\textbf{21.}] (1) \(\{f(x_n)\}\)为等比数列，公比\(-\mathrm e^{a\pi}\)；(2) 对\(a\geqslant\dfrac1{\sqrt{e^2-1}}\)，对一切\(n\in\mathbb N^*\)，有\(x_n<|f(x_n)|\)
\end{itemize}
